% ------------------------------------------------------------------------
% file `main-exercise-2-solution-body.tex'
%
%     solution of type `exercise' with id `2'
%
% generated by the `solution' environment of the
%   `xsim' package v0.21 (2022/02/12)
% from source `main' on 2026/11/14 on line 78
% ------------------------------------------------------------------------
\label {sol: subspaces and quotients of preordered vector spaces are preordered}We will prove each statement one-by-one.
\begin{enumerate}
    \item Since $W$ contains $0$ and $V^+$ contains $0$, we see that $W^+$ is not empty. If $x,y \in W^+$, then $x+y \in W$ and $x+y \in V^+$ since $x$ and $y$ individually belong to $V^+$. Finally, let $\lambda \in [0,\infty)$ and let $x \in W^+$. Then, $\lambda x \in W$ certainly and $\lambda x \in V^+$ so we're done.
    \item Since $V^+$ contains $0$, we see that $(V/W)^+$ contains $0$ as well since $\pi$ is a linear map, so it is non-empty. Next, let $x+W,y+W \in (V/W)^+$. We can choose $x$ and $y$ so that $x,y \in V^+$ and this means that $x+y \in V^+$. So, it follows that $(x+y)+W = (x+W) + (y+W) \in (V/W)^+$. Finally, let $\lambda \in [0,\infty)$ and let $x+W \in (V/W)^+$. We can choose $x \in V^+$ so $\lambda x \in V^+$ and this means that $\lambda x+W \in (V/W)^+$ and we're done.
\end{enumerate}
This completes the proof. In the case where $V$ is an ordered vector space, we proceed as follows:
\begin{enumerate}
    \item We need to show that:
    \[W^+ \cap (-W^+) = \{0\}.\]
    Only the inclusion from the left to the right needs to be shown, so let $x \in W^+ \cap (-W^+)$. Then, we see that $x \in V^+$ and $x \in -V^+$ so $x = 0$ and we're done.
    \item Let $x+W \in (V/W)^+ \cap (-(V/W)^+)$. Then, we can choose $x$ so that it comes from $V^+$. and the above implies that $x \in V^+ \cap (-V^+)$.
\end{enumerate}
This completes the whole proof.
